\documentclass{article} % For LaTeX2e
\usepackage[submission]{colm2026_conference}

\usepackage{microtype}
\usepackage{hyperref}
\usepackage{url}
\usepackage{booktabs}
\usepackage{amsmath}
\usepackage{amssymb}
\usepackage{enumitem}
\usepackage{graphicx}
\graphicspath{{./figs/}}
\usepackage{xcolor}
\usepackage[ruled,vlined,linesnumbered]{algorithm2e}
\usepackage{tikz}
\usetikzlibrary{arrows.meta,positioning}
\usepackage{lineno}

\definecolor{darkblue}{rgb}{0, 0, 0.5}
\hypersetup{colorlinks=true, citecolor=darkblue, linkcolor=darkblue, urlcolor=darkblue}

% algorithm2e sizing for narrow columns
\SetAlFnt{\small}
\SetAlCapFnt{\small}
\SetAlCapNameFnt{\small}

\title{Knowing but Not Saying: Probe-Saturation Gating and Engagement
  Schedules for the Knowing--Using Gap in Fine-Tuned Language Models}

% Authors are hidden automatically in submission mode (double-blind).
% For the camera-ready version, replace with the real author list.
\author{Anonymous Authors\\
  MLPR Research Team\\
  \texttt{mlpr@research.example} \\
}

\newcommand{\amemlf}{A_{\mathrm{mem}}^{\mathrm{lf}}}
\newcommand{\amemem}{A_{\mathrm{mem}}^{\mathrm{em}}}
\newcommand{\agen}{A_{\mathrm{gen}}}

\begin{document}

\ifcolmsubmission
\linenumbers
\fi

\maketitle

\begin{abstract}
When a language model is fine-tuned on new factual content, two capabilities
diverge: the model quickly becomes able to \emph{recall} an answer when the
question is posed in the training format (knowing), while its ability to
\emph{produce} the answer under distribution shift lags far behind (using).
We introduce MLPR, a lightweight monitoring-and-regularization layer that
makes this knowing--using gap measurable and actionable during training: a
linear probe at a mid layer $l^\star$ reads entity identities from hidden
states, and a saturation gate ramps in a probe-alignment loss
$\lambda(t)$ once a monitored signal crosses a threshold $\tau_0$. We
study Qwen2.5-0.5B and Qwen3-1.7B under an engagement-schedule grid of
ten fully instrumented runs on a synthetic entity-question corpus,
comparing two gate signals (output-likelihood recall vs.\
representation-space probe decodability) and three engagement schedules
(sustained, decayed, output-triggered). Three findings emerge. First,
the gate signal decides when engagement is possible: the probe signal
opens the gate at epoch 3--6 at both scales, 11--16 epochs before the
output signal crosses its own threshold in the replicates where it
ever does, and in one output-gated replicate it never crosses within
budget. Second, the engagement
schedule decides whether memorization completes: late output-triggered
engagement reaches SFT recall parity ($0.926$ vs.\ $0.934$; paired
McNemar $p{=}0.83$) and completes a previously budget-censored
regularized phase with 8 stable epochs at full probe weight, whereas
early sustained injection is high-variance ($0.109$--$0.805$) and
transient injection with decay to zero collapses recall to $0.012$
while the representation still decodes $94.8\%$ of entities---the
sharpest knowing--saying dissociation we observe. Third, free-
generation and generalization exact match stay at $\le 0.4\%$ in
\emph{every} arm at both scales even as teacher-forced recall reaches
$0.93$: the saying gap is scale-invariant, and all parameter updates
concentrate in upper-layer MLP projections. Code, telemetry, and
trained adapters are released.
\end{abstract}

\section{Introduction}
\label{sec:intro}

Fine-tuning is the standard route for injecting new knowledge into pretrained
language models, whether through supervised instruction tuning, domain
adaptation, or parameter-efficient updates \citep{hu2022lora}. A growing body
of evidence shows that this process memorizes far more readily than it
generalizes: models fit training-format facts quickly
\citep{tirumala2022memorization,carlini2023quantifying}, while behavioral
control of the acquired knowledge---answering paraphrases, chaining facts,
or refusing when unsure---arrives late or not at all
\citep{kandpal2023large,allen2023physics}. For practitioners this gap is
expensive: a model that \emph{knows} an entity (its answer is recoverable
from its internal states) but cannot \emph{say} it under mild distribution
shift will silently fail downstream, and standard observables cannot
distinguish the two regimes.

We propose to make this distinction measurable \emph{during} training, and
to act on it. Our approach, MLPR (Mid-Layer Probe Regularization), wraps a
parameter-efficient fine-tuning run with three coupled instruments. First,
a small linear probe $p_\phi$ attached to a mid layer $l^\star$ reads
entity identities from the backbone's hidden states; because the probe
trains on detached states during a warm-up phase, it is a pure observer
that cannot steer the backbone. Second, a teacher-forced likelihood metric
estimates recall saturation $\amemlf$ each epoch by scoring the answer
tokens given the training-format question---a ``knowing'' signal that, unlike
free-generation exact match, does not conflate knowledge with answer
extraction. Third, a saturation gate compares a monitored signal against a
threshold $\tau_0$ and, once crossed, ramps in a probe-alignment loss with
weight $\lambda(t)$ so that the backbone is regularized toward keeping acquired
identities linearly decodable at the mid layer. An eval-aware adaptive
controller completes the loop by damping the learning rate when held-out
loss degrades, subject to an explicit floor that keeps memorization
dynamics alive.

We instantiate MLPR on Qwen2.5-0.5B-Instruct and Qwen3-1.7B
\citep{qwen2024qwen25,qwen3}, training a grid of ten fully
instrumented runs (24 epochs; two 16-epoch censored replicates) on a
synthetic entity-question corpus with
a 1{,}600-entity vocabulary (per-step losses, per-epoch gate
evaluations, per-sample final evaluations, adapters and probe
checkpoints). The grid crosses two gate signals---output-likelihood
recall vs.\ representation-space probe decodability---with three
engagement schedules (sustained, decayed-to-zero, output-triggered),
completing the engagement-schedule arm that was budget-censored in our
earlier large-scale study. The resulting dataset of training dynamics
yields findings we summarize here and develop in
Section~\ref{sec:results}:

\begin{itemize}[leftmargin=1.4em,itemsep=1pt,topsep=2pt]
  \item \textbf{The gate signal decides when engagement is possible}
  (Section~\ref{sec:res:sat}). The probe's representation decodability
  crosses its threshold at epoch 3 (0.5B) and epochs 4--6 (1.7B), while
  the output-likelihood signal crosses $\tau_0{=}0.8$ only at epoch
  17--19---and in one output-gated replicate never crosses at all.
  Gating on the representation is scale-aware by construction; gating
  on the output is reachable at scale but unstable in crossing time.
  \item \textbf{The engagement schedule decides whether memorization
  completes} (Section~\ref{sec:res:sched}). Under identical data, seed,
  and optimizer, final recall spans $0.012$--$0.934$ at 1.7B. Late
  output-triggered engagement reaches SFT parity ($0.926$ vs.\ $0.934$,
  McNemar $p{=}0.83$) and completes the regularized phase---8 stable
  epochs at full probe weight---that our earlier large-scale run left
  budget-censored. Early sustained injection is high-variance
  ($0.109$, $0.805$); transient injection that decays to zero
  collapses recall to $0.012$ while the probe still decodes $94.8\%$
  of entities.
  \item \textbf{Knowing and saying dissociate at both scales}
  (Section~\ref{sec:res:gap}). Teacher-forced recall reaches
  $0.91$--$0.93$ and token-level recall $0.99$, yet free-generation EM
  stays $\le 0.012$ and generalization EM $\le 0.004$ in \emph{every}
  arm: the saying gap is invariant across a $3.4\times$ parameter
  range and across all schedules.
  \item \textbf{Probes detect memorization early and cheaply}
  (Section~\ref{sec:res:control}). The probe crosses $0.5$ decodability
  between epochs 4 and 7 everywhere, leading output recall by 8--18
  epochs (and leading arms whose output never crosses at all), with
  loss two orders of magnitude below chance.
  \item \textbf{Updates concentrate in upper-layer MLPs}
  (Section~\ref{sec:res:anatomy}). $\Delta W = \tfrac{\alpha}{r}BA$ per
  LoRA-adapted projection puts $65$--$71\%$ of update mass in MLP
  projections, key/value attention nearly untouched, consistent with
  factual storage accounts \citep{meng2022rome,geva2021ffn}.
\end{itemize}

Our contributions are: (i) a training-time instrumentation stack---probe,
likelihood-based recall metric, saturation gate with selectable signal,
and engagement schedules---that turns the knowing--using gap into an
observable, thresholdable quantity; (ii) a completed
engagement-schedule grid at two scales with paired McNemar analysis,
providing the first fully exercised regularized phase (8 epochs at
full probe weight) and its outcome parity with SFT; (iii) a
dissociation analysis separating representation decodability,
teacher-forced recall, free-generation EM, and chaining/intersection
transfer, including a representation-without-output regime induced by
transient engagement; and (iv) engineering lessons from instrumenting
real runs, including a documented failure mode and repair for
eval-aware learning-rate controllers (Appendix~\ref{app:spiral}).
All telemetry, code, and trained adapters are released.

\section{Related Work}
\label{sec:related}

\textbf{Probing internal representations.} Linear probes are a standard
instrument for asking what models encode
\citep{alain2017probes,belinkov2022probing}; control tasks guard against
probes that exploit surface statistics \citep{hewitt2019control}. Our
probe differs in role: trained \emph{alongside} the run on detached
states, its readings drive a control decision---activating a
regularization pathway---turning it from a diagnostic into a controller
input. Richer readouts \citep{bricken2023monosemanticity} could
substitute in future work.

\textbf{Knowledge storage and memorization dynamics.} Mechanistic and
editing studies localize factual knowledge in mid-to-upper MLP blocks
that behave as key--value memories
\citep{meng2022rome,meng2023memit,geva2021ffn,dai2022knowledgeneurons,
mitchell2022mend}; our anatomy analysis is consistent with this picture
from the direction of training dynamics. Synthetic entity corpora with
known ground truth have been used to study capacity and retrieval
\citep{allen2023physics}; our generator follows this tradition but
exposes per-epoch recall dynamics. Memorization of training data is
well documented \citep{carlini2023quantifying,zhang2017understanding},
large models memorize more before they overfit in loss terms
\citep{tirumala2022memorization}, and knowledge acquisition is strongly
scale-dependent \citep{kaplan2020scaling,kandpal2023large}; our grid
gives a fine-grained per-epoch view in the fine-tuning regime. The
dissociation we measure extends paraphrase-consistency findings
\citep{elazar2021measuring,mallen2023pararel} into training dynamics,
and recalls the competence--expression decoupling of grokking
\citep{power2022grokking}, here driven by data format rather than
weight decay.

\textbf{Adaptive training control.} Adaptive learning-rate schedules
respond to loss plateaus or gradient statistics; our controller adds a
held-out-loss criterion with a twist specific to the memorization
regime: because rising held-out loss is the very signal the gate waits
for, damping must be bounded lest the controller extinguish the
phenomenon it monitors. We are not aware of prior work analyzing this
controller--phenomenon interaction; we document a negative result (an
unbounded decay spiral) and its repair.

\section{Method: Mid-Layer Probe Regularization}
\label{sec:method}

MLPR augments parameter-efficient fine-tuning with an observation pathway
(a mid-layer probe and an online recall metric) and a control pathway (a
saturation gate and an eval-aware controller). Figure~\ref{fig:method}
shows the data flow; Algorithm~\ref{alg:mlpr} gives the training loop.

\begin{figure*}[t]
\centering
\resizebox{0.985\textwidth}{!}{%
% MLPR methodology overview — full-width figure* diagram.
% Custom TikZ icons (no external font/icon packages required):
%   snowflake = frozen, flame = trainable, gauge = saturation gate,
%   ramp inset = lambda(t), floor glyph = damping floor, slashed circle = detach.
% Shared by: fig_method.tex (standalone preview) and main paper via \input.
\newcommand{\mini}{\fontsize{5.8}{6.9}\selectfont}
\begin{tikzpicture}[
  x=1cm, y=1cm,
  fwd/.style={-{Stealth[length=5.5pt]}, semithick, black!72},
  tele/.style={-{Stealth[length=5pt]}, semithick, black!55, dashed},
  lamarr/.style={-{Stealth[length=6pt]}, semithick, gateC},
  blk/.style={draw=black!55, rounded corners=1.5pt, fill=gray!10,
              minimum width=0.68cm, minimum height=0.56cm,
              align=center, font=\scriptsize, inner sep=1pt},
  blkhl/.style={blk, fill=frozenBg, draw=frozenC, thick},
  chip/.style={draw=black!55, rounded corners=2pt, align=center,
               inner sep=2.5pt, font=\scriptsize},
  card/.style={chip, fill=white, text width=2.95cm},
  badge/.style={circle, fill=stepC, text=white, font=\tiny\bfseries,
                minimum size=0.34cm, inner sep=0pt},
  mini/.style={font=\mini},
]

% ================= palette =================
\definecolor{frozenC}{RGB}{68,119,170}   % blue  — frozen
\definecolor{frozenBg}{RGB}{232,239,247}
\definecolor{trainC}{RGB}{238,119,51}    % orange — trainable
\definecolor{gateC}{RGB}{170,51,119}     % magenta — gate / lambda
\definecolor{ctrlC}{RGB}{0,153,136}      % teal — controller
\definecolor{gradC}{RGB}{204,51,17}      % red — detach marker
\definecolor{openC}{RGB}{34,136,68}      % green — gate open
\definecolor{evC}{RGB}{238,170,0}        % amber — evidence card
\definecolor{stepC}{RGB}{85,85,85}
\definecolor{anchorBg}{RGB}{253,231,214}

% ================= icons =================
\tikzset{
  flake/.pic={
    \foreach \a in {0,60,120}{\draw[frozenC, line width=0.55pt] (\a-180:0.145) -- (\a:0.145);}
    \foreach \a in {0,60,...,300}{
      \draw[frozenC, line width=0.45pt] (\a:0.085) -- ++(\a+42:0.055);
      \draw[frozenC, line width=0.45pt] (\a:0.085) -- ++(\a-42:0.055);}
    \fill[frozenC] (0,0) circle(0.03);
  },
  flame/.pic={
    \fill[trainC] (0,0.19) .. controls (-0.15,0.04) and (-0.13,-0.07) .. (-0.02,-0.17)
        .. controls (0.10,-0.08) and (0.13,0.02) .. (0,0.19) -- cycle;
    \fill[white] (0,0.055) .. controls (-0.055,-0.02) and (-0.055,-0.075) .. (0,-0.125)
        .. controls (0.055,-0.07) and (0.055,-0.02) .. (0,0.055) -- cycle;
  },
  gauge/.pic={
    \draw[black!55, line width=0.5pt] (-0.30,0) arc[start angle=180, end angle=0, radius=0.30];
    \draw[black!55, line width=0.5pt] (-0.30,0) -- (0.30,0);
    \foreach \ang in {180,135,90,45,0}{\draw[black!45, line width=0.4pt] (\ang:0.26) -- (\ang:0.30);}
    \draw[gateC, line width=1.1pt] (36:0.21) -- (36:0.335);   % tau_0 = 0.8
    \draw[black!70, line width=1.2pt] (0,0) -- (34.9:0.245);  % needle at 0.805
    \fill[black!60] (0,0) circle(0.032);
  },
  ramp/.pic={
    \draw[->, black!55, line width=0.5pt] (-0.05,0) -- (1.42,0);
    \draw[->, black!55, line width=0.5pt] (0,-0.06) -- (0,0.86);
    \draw[gateC, line width=1.1pt] (0,0) -- (0.56,0) -- (1.36,0.76);
    \fill[gateC] (0.56,0) circle(0.035);
    \node[mini, black!60, below] at (0.56,-0.03) {$t_0$};
    \node[mini, black!60, left]  at (-0.03,0.74) {$\lambda_0$};
  },
  floorplt/.pic={
    \draw[->, black!55, line width=0.5pt] (-0.05,0.72) -- (-0.05,0);
    \draw[->, black!55, line width=0.5pt] (-0.05,0) -- (1.52,0);
    \draw[ctrlC, line width=1.1pt] (0,0.62) -- (0.30,0.40) -- (0.58,0.27) -- (0.82,0.21)
        -- (1.05,0.19) -- (1.45,0.19);
    \draw[black!45, dashed, line width=0.5pt] (-0.05,0.19) -- (1.45,0.19);
    \node[mini, black!60, below] at (1.18,0.0) {$m_{\min}$};
    \node[mini, black!60, left]  at (-0.05,0.62) {$m$};
  },
  det/.pic={ % detach: slashed circle
    \fill[white] (0,0) circle(0.085);
    \draw[gradC, line width=0.8pt] (0,0) circle(0.085);
    \draw[gradC, line width=0.8pt] (135:0.115) -- (-45:0.115);
  },
}

% ================= band backgrounds =================
\fill[blue!3, rounded corners=4pt]   (0.10,-1.50) rectangle (16.98,3.48);
\fill[black!3, rounded corners=4pt]  (0.10,-6.02) rectangle (16.98,-1.72);
\node[anchor=west, font=\scriptsize\bfseries, black!55] at (0.32,3.26)
  {DATA \& OBSERVATION PATHWAY \; (per batch)};
\node[anchor=west, font=\scriptsize\bfseries, black!55] at (0.32,-1.78)
  {CONTROL PATHWAY \; (per epoch)};

% ================= worked-example query card =================
\node[card, anchor=north west] at (0.35,2.62) {%
  {\scriptsize\bfseries Worked example}\\[2pt]
  {\mini\itshape Q:} Which regulatory authority oversees
  \colorbox{anchorBg}{\mini Institution\_Alpha\_175}?\\[3pt]
  {\mini\itshape A:} \texttt{Answer:}\\[0.5pt]
  {\mini\hspace*{1.05em} Regulator\_Beta\_044}\\[3pt]
  {\mini probe label: entity id 243}\\
  {\mini anchor token: chars 36--57}};

% ================= backbone chain =================
\node[circle, draw=black!55, fill=gray!8, minimum size=0.42cm,
      font=\scriptsize, inner sep=0pt] (emb) at (4.12,1.9) {$x$};
\node[blk] (L1)  at (4.94,1.9) {$1$};
\node[blk] (L2)  at (5.78,1.9) {$2$};
\node[font=\scriptsize, black!55] at (6.42,1.9) {$\cdots$};
\node[blkhl] (LS) at (7.06,1.9) {$l^\star$};
\node[blk] (LL)  at (8.34,1.9) {$L$};
\pic at (4.98,2.80) {flake};
\node[anchor=west, font=\mini, black!60] at (5.18,2.80)
  {decoder layers $1\ldots L$ --- backbone $\theta$ \textbf{frozen}};

% LoRA branch (trainable)
\node[chip, fill=trainC!10, draw=trainC!80!black, minimum width=1.02cm] (Abox) at (5.00,-0.90)
  {\scriptsize LoRA $A$\\{\mini $d\!\to\!r{=}16$}};
\node[chip, fill=trainC!10, draw=trainC!80!black, minimum width=1.02cm] (Bbox) at (6.42,-0.90)
  {\scriptsize LoRA $B$\\{\mini $r\!\to\!d$}};
\pic at (4.66,-0.615) {flame};
\pic at (6.08,-0.615) {flame};
\node[font=\mini, black!60, align=center, text width=3.0cm] at (5.62,-0.22)
  {$\Delta W{=}\alpha r^{-1}BA$ per layer\\(7 projections: attn + MLP)};
\draw[fwd] (emb.south) -- (4.12,-0.90) -- (Abox.west);
\draw[fwd] (Abox.east) -- (Bbox.west);
\node[circle, draw=black!60, fill=white, inner sep=0.4pt, font=\tiny] (plus) at (7.62,1.9) {$+$};
\draw[fwd] (Bbox.east) -- (7.62,-0.90) -- (plus.south);

% LM head / answer / L_ce
\node[chip, fill=gray!8, minimum width=0.92cm] (head) at (9.32,1.9) {\scriptsize LM head};
\node[chip, fill=white, minimum width=2.42cm] (ans) at (11.55,1.9)
  {\scriptsize Answer: Regulator\_Beta\_044};
\node[chip, fill=blue!8, minimum width=1.12cm, minimum height=0.72cm] (lce) at (14.00,1.9)
  {\scriptsize $\mathcal{L}_{\mathrm{ce}}$\\{\mini SFT loss}};
\draw[fwd] (emb.east) -- (L1.west);
\draw[fwd] (LL.east) -- (head.west);
\draw[fwd] (head.east) -- (ans.west);
\draw[fwd] (ans.east) -- (lce.west);

% epoch-end eval node
\node[chip, fill=gray!6, text width=2.0cm] (evaln) at (15.75,2.94)
  {\scriptsize\bfseries epoch-end eval\\{\mini 256-prompt recall subsample}\\
   {\mini $A_{\mathrm{mem}}^{\mathrm{lf}},\,A_{\mathrm{mem}}^{\mathrm{em}},\,A_{\mathrm{gen}}$}};
\draw[tele] (ans.north) |- (evaln.west);

% probe pathway
\node[chip, fill=trainC!10, draw=trainC!80!black, text width=2.52cm] (probe) at (6.05,0.50)
  {\scriptsize\bfseries linear probe $p_\phi$\\{\mini $d\!\to\!|V_e|{=}1600$ · LR $2{\times}10^{-3}$}};
\pic at (4.62,0.82) {flame};
\node[chip, fill=white, minimum width=2.26cm] (pred) at (9.00,0.50)
  {\scriptsize entity: Regulator\_Beta\_044\\{\mini probe argmax = id 243}};
\node[chip, fill=trainC!10, draw=trainC!80!black, minimum width=1.30cm,
      minimum height=0.72cm] (lpr) at (11.42,0.50)
  {\scriptsize $\mathcal{L}_{\mathrm{probe}}$\\{\mini entity CE}};
\draw[fwd] (LS.south) -- (7.06,0.84);
\node[mini, black!60, anchor=east] at (6.94,1.12) {$h_{l^\star}$ (anchor)};
\pic at (7.06,1.42) {det};
\node[mini, gradC, anchor=east] at (6.90,1.42) {detached (warm-up)};
\draw[fwd] (probe.east) -- (pred.west);
\draw[fwd] (pred.east) -- (lpr.west);

% ================= control pathway =================
% gate
\node[chip, fill=gateC!8, draw=gateC, thick, anchor=west, text width=2.42cm,
      minimum height=1.24cm] (gate) at (0.30,-2.70)
  {\scriptsize\bfseries saturation gate\\[2pt]
   {\mini $\rho{\in}\{A_{\mathrm{mem}}^{\mathrm{lf}},\,\rho_{\mathrm{rep}}\}$}\\[1pt]
   {\mini $\ge\ \tau_0$\;?}};
\pic at (3.30,-2.62) {gauge};

% closed / open chips
\node[chip, fill=gray!8, minimum width=1.52cm] (closed) at (0.95,-3.74)
  {\mini closed: $\lambda{=}0$\\{\mini pure SFT}};
\node[chip, fill=openC!10, draw=openC!80!black, minimum width=1.56cm] (open) at (2.72,-3.74)
  {\mini open at $t_0$:\\{\mini ramp $\lambda(t)$}};
\draw[fwd] (gate.south) -- ++(0,-0.16) -| (closed.north);
\draw[fwd] (gate.south) -- ++(0,-0.16) -| (open.north);

% evidence card
\node[chip, fill=evC!12, draw=evC!75!black, text width=3.42cm] (ev) at (1.95,-4.94)
  {\scriptsize\bfseries worked example --- 1.7B run\\[1.5pt]
   {\mini ep 17: $A_{\mathrm{mem}}^{\mathrm{lf}}{=}0.885 \ge \tau_0$}\\[0.5pt]
   {\mini gate opens at $\lambda_0{=}0.3$}\\[0.5pt]
   {\mini ep 18: signal re-enters band}\\[0.5pt]
   {\mini $\lambda$: $0.3{\to}0.075{\to}0.3$ (live)}};

% lambda ramp inset
\pic at (4.30,-2.62) {ramp};
\node[font=\mini, black!60, align=center, anchor=north, text width=2.6cm] at (4.45,-2.92)
  {$\lambda(t){=}\lambda_0\,\min\!\bigl(1,\tfrac{\rho(t)-\tau_0}{\Delta_0}\bigr)$\\[1pt]
   {\mini $\lambda_0{=}0.3$; decay variant:}\\[0.5pt]
   {\mini $\times\max(0,1{-}\tfrac{t-t_0}{D})$}};

% total loss
\node[chip, fill=yellow!14, draw=black!60, thick, minimum width=2.56cm,
      minimum height=1.18cm] (loss) at (8.15,-2.62)
  {\scriptsize\bfseries total objective\\[1.5pt]
   {\scriptsize $\mathcal{L}=\mathcal{L}_{\mathrm{ce}}+\lambda(t)\,\mathcal{L}_{\mathrm{probe}}$}};
\draw[lamarr] (5.80,-2.62) -- (6.87,-2.62);
\node[badge] at (6.33,-2.62) {5};
\node[mini, gateC, anchor=north] at (6.33,-2.84) {$\lambda(t)$};

% controller + optimizer
\node[chip, fill=ctrlC!8, draw=ctrlC, thick, text width=2.92cm] (ctrl) at (12.50,-2.62)
  {\scriptsize\bfseries eval-aware controller\\[2pt]
   {\mini boost LR while held-out improves;}\\[0.5pt]
   {\mini decay $\times\,0.7$ after degradation;}\\[0.5pt]
   {\mini floor $m\leftarrow\max(0.5,\,m\gamma)$}};
\pic at (12.50,-4.18) {floorplt};
\node[circle, draw=black!60, fill=gray!10, minimum size=1.06cm, font=\tiny, inner sep=0pt]
  (opt) at (15.42,-2.62) {\textbf{AdamW}};
\draw[fwd] (loss.east) -- node[above, mini, black!60] {gradients} (ctrl.west);
\draw[fwd] (ctrl.east) -- node[above, mini, black!60] {LR $\times\,m(t)$} (opt.west);

% trainable scope chip
\node[chip, fill=trainC!8, draw=trainC!80!black, text width=1.9cm] (scope) at (15.42,-3.95)
  {\mini \tikz{\pic[scale=0.62]{flame}}\;updates only $A,B,\phi$\\[1pt]
   \tikz{\pic[scale=0.62]{flake}}\;$\theta$ stays frozen};
\draw[fwd] (opt.south) -- (scope.north);

% ================= telemetry / loss feeds =================
\draw[tele] (10.60,1.64) -- (10.60,-1.55) -| (1.60,-1.97);
\node[badge] at (10.60,-0.36) {4};
\draw[fwd] (lce.south) -- (14.00,-1.32) -- (8.72,-1.32) -- (8.72,-2.03);
\draw[fwd] (lpr.south) -- (11.42,-1.14) -- (7.62,-1.14) -- (7.62,-2.03);

% ================= legend =================
\node[chip, fill=white, text width=2.82cm, anchor=north] (leg) at (15.42,-4.50) {%
  {\mini\bfseries legend}\\[1.5pt]
  {\mini \tikz{\pic[scale=0.62]{flake}} frozen $\theta$ \;\; \tikz{\pic[scale=0.62]{flame}} trainable}\\[1pt]
  {\mini \tikz{\draw[black!55,semithick,dashed,-{Stealth[length=3pt]}](0,0)--(0.42,0);} eval-only signal}\\[1pt]
  {\mini \tikz{\pic[scale=0.62]{det}} detached (no grad)}};

% ================= walkthrough badges =================
\node[badge] at (3.70,1.9) {1};
\node[badge] at (7.06,2.44) {2};
\node[badge] at (7.06,1.10) {3};
\node[badge] at (15.42,-3.36) {6};

\end{tikzpicture}
}
\caption{MLPR methodology with a worked example.
\textbf{Top --- data \& observation pathway (per batch):} a memorization
sample from $\mathcal{D}_{\mathrm{mem}}$ (left card; anchor token
highlighted) is processed by the frozen backbone (snowflake: $\theta$
receives no gradient) augmented with trainable low-rank adapters
(flame: $A,B$, $r{=}16$, applied to all seven attention and MLP
projections, merged at $\oplus$). The mid-layer state $h_{l^\star}$ at
the anchor token is read by a linear probe $p_\phi$ (flame) that
predicts the entity identity; during warm-up the probe observes a
\emph{detached} copy of the state ($\oslash$, no gradient into
$\theta$). \textbf{Bottom --- control pathway (per epoch):} the
saturation gate (gauge) compares teacher-forced recall
$\rho=\max(\amemlf,\amemem)$ against $\tau_0{=}0.8$; once it opens,
$\lambda(t)$ ramps the probe-alignment loss into the total objective
(ramp inset), and an eval-aware controller modulates the learning rate
under a damping floor $m\ge m_{\min}{=}0.5$ (floor inset). Numbered
badges trace the example: (1)~sample enters; (2)~frozen+LoRA forward;
(3)~detached mid-layer probe read; (4)~epoch-end recall evaluation
(208-prompt subsample); (5)~gated loss weight; (6)~optimizer updates
trainable parameters only. The amber card reports the realized 7B run:
the gate opened at epoch 24 ($\amemlf{=}0.805\ge\tau_0$) and $\lambda$
ramped $0.015\!\to\!0.075\!\to\!0.225$ toward $\lambda_0{=}0.3$,
whereas the 0.5B run never crossed the threshold ($\max 0.39$).}
\label{fig:method}
\end{figure*}

\subsection{Setup and objective}
\label{sec:method:setup}

Let $f_\theta$ be a pretrained decoder-only LM, and let LoRA
\citep{hu2022lora} add low-rank updates
$\Delta W = \tfrac{\alpha}{r} BA$ to selected weight matrices, with $A,B$
trainable and $\theta$ frozen. The supervised objective on the
memorization corpus $\mathcal{D}_{\mathrm{mem}}$ of question--answer
pairs is the token-level cross-entropy
\begin{equation}
  \mathcal{L}_{\mathrm{ce}}
  = -\frac{1}{|B|}\sum_{(x,y)\in B}
    \log p_\theta\!\bigl(y \mid x\bigr),
  \label{eq:ce}
\end{equation}
where $x$ is the training-format question and $y$ the answer with
\texttt{Answer:} suffix. At each optimizer step the total loss is
\begin{equation}
  \mathcal{L} = \mathcal{L}_{\mathrm{ce}}
  + \lambda(t)\, \mathcal{L}_{\mathrm{probe}}\bigl(p_\phi, f_\theta\bigr),
  \label{eq:total}
\end{equation}
where $\lambda(t)\!\in\![0,\lambda_0]$ is the gate weight
(Section~\ref{sec:method:gate}) and $\mathcal{L}_{\mathrm{probe}}$ is a
cross-entropy between the probe's entity prediction and the ground-truth
entity id of the sample in the batch, computed on the mid-layer hidden
state $h_{l^\star}$ \emph{with gradients flowing into the backbone} once
$\lambda>0$ (before activation the probe receives no gradient path into
$\theta$, cf.\ Eq.~\ref{eq:total} with $\lambda{=}0$).

\subsection{Probe warm-up and the likelihood recall metric}
\label{sec:method:probe}

The probe $p_\phi:\mathbb{R}^{d}\!\to\!\mathbb{R}^{V_e}$ is a single
linear layer over the entity vocabulary $V_e$ (here $|V_e|{=}1600$),
trained from the first step on detached states at
$l^\star{\approx}L/2$ with its own learning rate
($10^{-3}$ vs.\ $2{\times}10^{-5}$ for the backbone), so that it
converges quickly enough to be informative early. Because the probe
gradient is detached during warm-up, it is a pure observer; only after
gate activation does $\mathcal{L}_{\mathrm{probe}}$ back-propagate into
$\theta$, at which point it encourages the backbone to keep acquired
entity identities linearly decodable at $l^\star$.

Free-generation exact match is a poor saturation signal: a model may
encode the correct answer yet fail to emit it in the right surface form.
We therefore measure recall with a teacher-forced likelihood probe: for
each memorization prompt (answer suffix stripped), the answer tokens are
scored in a single forward pass and the greedy per-sample match is
\begin{equation}
  \mathrm{lf}(x, y) = \mathbb{1}\bigl[
    \arg\max_w p_\theta(w \mid x) = y_1 \bigr],
  \label{eq:lf}
\end{equation}
with a token-level variant that credits partial answers. The epoch-level
recall is $\amemlf = \mathbb{E}_{\mathcal{S}}[\mathrm{lf}]$ on a fixed
subsample $\mathcal{S}$ of $\mathcal{D}_{\mathrm{mem}}$ (208 prompts).
We additionally log the free-generation exact match
$\amemem$, substring containment, paraphrased generalization accuracy
$\agen$ on $\mathcal{D}_{\mathrm{gen}}$, and the probe's own
train/held-out accuracies, giving a five-channel view of the same
underlying state change.

\subsection{Saturation gate and $\lambda(t)$}
\label{sec:method:gate}

Let $t$ index the training step and $T$ the total step budget. The gate
uses the maximum of the likelihood and free-generation recall signals,
$\rho(t) = \max(\amemlf, \amemem)$, and a threshold $\tau_0$. Before
$\rho \ge \tau_0$ the gate is closed ($\lambda{=}0$) and training is
pure SFT. After the first crossing at step $t_0$, the weight ramps
linearly over a fraction $\delta_0$ of the budget,
\begin{equation}
  \lambda(t) = \lambda_0 \cdot
    \min\!\Bigl(1,\ \tfrac{t - t_0}{\delta_0 T}\Bigr),
  \qquad t \ge t_0,
  \label{eq:ramp}
\end{equation}
so the regularization is introduced gradually ($\lambda_0{=}0.3$,
$\tau_0{=}0.8$, $\delta_0{=}0.1$ here). The maximum over the two recall
modes makes the gate robust to the answer-extraction artifact
\citep{mallen2023pararel}: in pilot runs the free-generation signal
plateaued at $0$--$12\%$ and never opened a gate that required it, while
the likelihood signal crossed the same threshold at scale
(Section~\ref{sec:res:sat}).

\subsection{Eval-aware adaptive control with a damping floor}
\label{sec:method:control}

Long memorization runs exhibit a characteristic trajectory: training
loss plateaus after the format is learned, and held-out loss---bottoming
early---then rises monotonically as capacity is redirected to training
entities \citep{tirumala2022memorization}. An adaptive controller that
treats the plateau as under-optimization may boost the learning rate
exactly while validation degrades, and one that treats the rise as
divergence may decay it without bound. Both failure modes occurred in
our pilot runs. The controller we ship therefore (i) boosts only while
held-out loss is not degrading, (ii) decays by a fixed factor after
repeated degradation episodes, and (iii) enforces a \emph{damping floor}:
the cumulative LR multiplier $m(t)$ applied by eval-driven decays is
clamped to $m(t) \ge m_{\min}$ with $m_{\min}{=}0.5$,
\begin{equation}
  m(t) \leftarrow \max\!\bigl(m_{\min},\ m(t^-)\cdot \gamma\bigr),
  \label{eq:floor}
\end{equation}
with decay factor $\gamma{=}0.7$. The floor is the load-bearing fix: in
the memorization regime, held-out divergence is the phenomenon under
study, not a fault, and unbounded damping freezes it
(Section~\ref{sec:res:control} quantifies the spiral we observed without
the floor). Every decision is logged per epoch, and all controller
actions are failure-isolated so that instrumentation cannot crash
training.

\begin{algorithm}[t]
\caption{MLPR training loop (one epoch)}
\label{alg:mlpr}
\KwIn{corpora; LoRA backbone $f_\theta$; probe $p_\phi$; gate
  \eqref{eq:ramp}; controller with floor \eqref{eq:floor}}
probe warm-up on detached $h_{l^\star}$\;
\ForEach{batch $B \subset \mathcal{D}_{\mathrm{mem}}$}{
  $\mathcal{L} \leftarrow \mathcal{L}_{\mathrm{ce}}
    + \lambda(t)\,\mathcal{L}_{\mathrm{probe}}$; update $\{A,B,\phi\}$\;
}
$\amemlf \leftarrow$ \eqref{eq:lf} on subsample; evaluate
$\amemem, \agen$\;
open gate if $\rho \ge \tau_0$; clamp controller $m \ge m_{\min}$; log\;
\end{algorithm}

\section{Experimental Setup}
\label{sec:setup}

\textbf{Corpus.} We generate a synthetic entity-question corpus in the
spirit of controlled knowledge-store studies \citep{allen2023physics}: a
vocabulary of $1{,}600$ invented entities, each assigned to generated
questions with deterministic single-token-friendly answers. The
memorization split $\mathcal{D}_{\mathrm{mem}}$ contains 1{,}200
examples in a fixed template with an \texttt{Answer:} suffix; the
generalization split $\mathcal{D}_{\mathrm{gen}}$ (450 prompts) asks the
same facts through paraphrase and chaining templates never seen in
training, so exact match on it measures \emph{use} rather than recall.
Dataset construction details and templates are in
Appendix~\ref{app:dataset}.

\textbf{Models and training.} We train Qwen2.5-0.5B-Instruct and
Qwen2.5-7B-Instruct \citep{qwen2024qwen25} with LoRA
($r{=}16$, $\alpha{=}32$, dropout $0.05$) on all attention and MLP
projections, probe layer $l^\star{=}12$ and $l^\star{=}14$
respectively, backbone LR $2{\times}10^{-5}$ (cosine, $10\%$ warmup),
effective batch 16, probe LR $10^{-3}$, gate parameters
$\tau_0{=}0.8$, $\delta_0{=}0.1$, $\lambda_0{=}0.3$, and the controller
of Section~\ref{sec:method:control} with $m_{\min}{=}0.5$. Each model
trains for up to 50 epochs (75 steps/epoch) on a single H100 in bf16;
the 0.5B run completed all 50 epochs (3{,}750 steps), while the 7B run
was terminated by compute-budget exhaustion during epoch 28, after 27
completed epochs with the gate open. All telemetry---per-epoch gate
evaluations, controller decisions, probe states, matrices, and the
complete stdout logs---is persisted to the Hugging Face Hub
(Appendix~\ref{app:artifacts}).

\textbf{Metrics.} Each epoch we log: likelihood recall $\amemlf$
\eqref{eq:lf}; token-level recall; free-generation exact match
$\amemem$; substring containment of the target in the generation;
paraphrase exact match $\agen$; the probe's accuracy on memorized vs.\
held-out entity states; mean training loss; and held-out eval loss.
Table~\ref{tab:main} summarizes the endpoints of both runs.

\begin{table}[t]
\centering
\small
\begin{tabular*}{\columnwidth}{@{\extracolsep{\fill}} l cc}
\toprule
 & 0.5B (complete) & 7B (ep.\ 27) \\
\midrule
epochs completed & 50 & 27 \\
recall $\amemlf$ final & 0.380 & 0.875 \\
token recall final & 0.900 & 0.987 \\
free-gen EM $\amemem$ final & 0.025 & 0.020 \\
substring containment & 0.280 & 0.765 \\
paraphrase EM $\agen$ & 0.000 & 0.000 \\
gate opened (epoch) & never & 24 \\
$\lambda(t)$ final & 0.000 & 0.225 \\
eval loss (min $\to$ final) & $3.36 \to 4.64$ & $3.21 \to 4.41$ \\
LR multiplier final & 0.50 & 0.50 \\
\bottomrule
\end{tabular*}
\caption{Run endpoints. The 0.5B model never saturates recall within
budget; the 7B model crosses $\tau_0{=}0.8$ at epoch 24 and the gate
opens. Free-generation exact match stays near zero at both scales even
as likelihood recall reaches $0.875$; at 7B the generated answers
contain the target as a substring $76.5\%$ of the time, isolating the
failure in answer assembly.}
\label{tab:main}
\end{table}

\section{Results}
\label{sec:results}

\subsection{Recall saturation is scale-dependent, and the gate opens at scale}
\label{sec:res:sat}

Figure~\ref{fig:sat}a overlays the two runs' likelihood recall. The 7B
curve is S-shaped: a slow phase through epoch 11 ($\amemlf \le 0.11$),
a rapid phase through epoch 22 ($0.77$), and crossing
$\tau_0{=}0.8$ at epoch 24 ($0.805$), then $0.825$ and $0.875$ at
epochs 25 and 27. The 0.5B curve rises roughly linearly after a long
delay ($0.055$ at epoch 16, $0.22$ at epoch 31, $0.39$ at epoch 50,
still improving by $\sim$0.01 per 5 epochs) and never reaches the
threshold within twice the 7B crossing budget---a per-epoch view of the
scale dependence of knowledge acquisition reported at pretraining scale
\citep{kaplan2020scaling,kandpal2023large,allen2023physics}, within the
\emph{same} architecture family, corpus, and optimizer.

\begin{figure*}[t]
\centering
\includegraphics[width=\textwidth]{fig_saturation.pdf}
\caption{Scale-dependent saturation and live gate activation.
\textbf{(a)} Likelihood recall $\amemlf$ vs.\ epoch for both scales;
the 7B model crosses $\tau_0{=}0.8$ (dashed) at epoch 24 (star), the
0.5B model plateaus at $0.39$ after 50 epochs.
\textbf{(b)} The 7B knowing--saying dissociation: teacher-forced recall
climbs to $0.875$ while paraphrase exact match ($\agen$) remains at
$0.000$ and free-generation EM stays $\le 0.07$.
\textbf{(c)} The gate weight $\lambda(t)$: closed for the entire 0.5B
run; opened at epoch 24 at 7B with the ramp
$0.015 \to 0.075 \to 0.225$ observed over epochs 24--27 (the run was
budget-truncated before the ramp completed to $\lambda_0{=}0.3$).}
\label{fig:sat}
\end{figure*}

Because the crossing happened mid-run, the activation transient was
logged live: the gate recorded
$(\amemlf, \lambda) = (0.805, 0.015)$ at epoch 24,
$(0.825, 0.075)$ at epoch 25, and $(0.875, 0.225)$ at epoch 27
(Fig.~\ref{fig:sat}c). To our knowledge this is an unusually direct
capture of a ``memorization threshold'' crossing during a real
fine-tuning run; the ramp toward $\lambda_0{=}0.3$ was cut short by the
compute budget, so the regularized phase (probe loss flowing into the
backbone) is exercised but not yet studied at length---we return to this
in Section~\ref{sec:res:post}.

\subsection{Knowing is not saying}
\label{sec:res:gap}

Figure~\ref{fig:gap} plots the five recall/production channels for both
scales. Two dissociations are visible. First, \emph{knowing vs.\
saying}: at 7B the teacher-forced recall reaches $0.875$ while
free-generation exact match stays $\le 0.07$ and paraphrase exact match
is exactly $0.000$ at every epoch; the 0.5B model shows the same
qualitative pattern ($0.39$ vs.\ $0.025$ vs.\ $0.000$). Second,
\emph{storage vs.\ assembly}: at epoch 27 the 7B generations contain
the target string as a substring $76.5\%$ of the time while exact match
is $2\%$---the entity identity is present in the output but framed
wrong (extra tokens, wrong surface form), which is precisely the
failure mode that free-generation EM conflates with ignorance and that
the likelihood metric separates. Token-level recall ($0.987$ at 7B)
confirms that the first answer token---the identity-bearing
token---is almost always correct under forcing; what fails is the
assembly of a well-formed answer under free decoding.

\begin{figure}[t]
\centering
\includegraphics[width=0.80\columnwidth]{fig_knowing_saying.pdf}
\caption{Knowing vs.\ saying at both scales. Shaded region: the
knowing--using gap between teacher-forced recall and free-generation
production. Paraphrase exact match ($\agen$, dashed) never leaves zero
at either scale.}
\label{fig:gap}
\end{figure}

The knowing--using gap is thus not a delay that more epochs alone close
at small scale: the 0.5B model's gap persists through 50 epochs
($0.38$ vs.\ $0.025$), and the 7B model's gap is $0.855$ in recall
terms at the moment the gate opens. Regularization or decoding
interventions that target the \emph{saying} pathway (e.g., the
$\lambda$-gated probe loss, or answer-assembly training) are therefore
better placed than additional memorization pressure.

\subsection{Dynamics: eval reversal and bounded adaptive control}
\label{sec:res:control}

Figure~\ref{fig:dyn} shows the shared training dynamics. Held-out eval
loss bottoms at epoch 3 for both scales ($3.36$ and $3.21$) and then
rises monotonically ($4.64$ and $4.41$ by the last completed epoch)---
the memorization signature. Training loss falls from $\sim$13 to
$\sim$0.9 (0.5B) and $\sim$0.85 (7B), with the 7B run reaching lower
CE earlier ($\sim$0.30 mean by epoch 25).

\begin{figure*}[t]
\centering
\includegraphics[width=\textwidth]{fig_dynamics.pdf}
\caption{Training dynamics. \textbf{(a)} Mean per-epoch training loss.
\textbf{(b)} Held-out eval loss reverses at epoch 3 at both scales and
rises thereafter---the memorization signature the gate waits for.
\textbf{(c)} Controller behavior: both runs decay the LR twice under
held-out divergence, then hold at the $0.5\times$ damping floor
(dashed).}
\label{fig:dyn}
\end{figure*}

The controller behaved identically at both scales: two eval-driven
decays (epochs 6 and 8; multiplier $1.0 \to 0.7 \to 0.5$) followed by
a hold at the floor for the rest of training, with plateau boosts
correctly suppressed while eval loss degraded. This is the repaired
behavior. Without the floor (pilot runs), the same controller applied
consecutive $0.7\times$ decays and drove the multiplier to $0.17$,
freezing memorization with $\amemlf$ pinned at $\sim$0.02 and making
$\tau_0$ unreachable; Appendix~\ref{app:spiral} documents the failure
and the validation sequence (unit tests plus a dedicated 5-epoch GPU
gate check) that the repair passed before these production runs.

\subsection{Probes read memorization before the model can say it}
\label{sec:res:probe}

Figure~\ref{fig:probe}a shows the probe's epoch-end accuracy on
memorized vs.\ held-out entity states. At both scales the probe
separates the two populations within a few epochs ($7$--$13\%$ vs.\
$\approx 0\%$ at chance $1/1600$), and the separation tracks the
saturation curve: at 7B the probe's memorization accuracy roughly
doubles after the gate opens. The probe loss
(Fig.~\ref{fig:probe}b) falls from chance $\ln 1600 \approx 7.38$ to
below chance by orders of magnitude---reaching $\sim$$10^{-3}$ with
$100\%$ train-batch accuracy at 7B---while remaining a pure observer
during warm-up. Two practical readings follow: the probe is a cheap
early-warning monitor (its signal is unambiguous epochs before any
behavioral metric moves), and its health is verifiable online (loss
far below chance, held-out accuracy at chance), addressing the classic
probe-reliability concern \citep{hewitt2019control}.

\begin{figure}[t]
\centering
\includegraphics[width=0.82\columnwidth]{fig_probe.pdf}
\caption{Probe health. \textbf{(a)} Epoch-end probe accuracy on
memorized ($\mathcal{D}_{\mathrm{mem}}$ subsample) vs.\ held-out
($\mathcal{D}_{\mathrm{gen}}$) entity states; chance is $1/1600$.
\textbf{(b)} Probe loss (log scale) falls far below chance at both
scales during detached warm-up.}
\label{fig:probe}
\end{figure}

\subsection{Update anatomy: upper-layer MLPs carry the memorization}
\label{sec:res:anatomy}

From the released adapters we compute the effective weight update
$\Delta W = \tfrac{\alpha}{r} BA$ per adapted projection
(Fig.~\ref{fig:anat}). The pattern is consistent across scales and
consistent with the storage literature: MLP projections carry
$69.7\%$ (0.5B) and $73.3\%$ (7B) of total update mass
$\sum_l \lVert \Delta W \rVert_F$; the upper half of layers carries
$57.8\%$ and $56.1\%$; and the top three modules are
gate-, up-, and down-projection at both scales, with key/value
attention projections nearly untouched. The 7B adapter analyzed here
is the epoch-26 state, i.e., one to two epochs \emph{after} gate
activation, so the concentration is not an artifact of an
un-activated regime.

\begin{figure}[t]
\centering
\includegraphics[width=0.78\columnwidth]{fig_anatomy.pdf}
\caption{Update anatomy $\lVert \Delta W \rVert_F$ per
(layer, module). \textbf{(a)} 0.5B final adapter (epoch 50);
\textbf{(b)} 7B adapter at epoch 26 (post-activation). MLP projections
(gate/up/down) dominate at both scales; color scales are per-panel.}
\label{fig:anat}
\end{figure}

\subsection{The regularized phase and its truncation}
\label{sec:res:post}

Honesty about scope: the gate opened at 7B only $\sim$3 epochs before
the compute budget terminated the run, so the post-activation phase
(probe loss at $\lambda$ up to $0.225$ flowing into the backbone) is
\emph{exercised and logged} but not yet studied over the tens of
epochs its design assumes. Within the observed window the phase is
stable (no loss or gradient spikes; $\amemlf$ continues to rise;
no further controller interventions). Whether the probe-regularized
phase closes the \emph{saying} gap---the motivation for the
design---remains the central open question for the completed 7B run;
the 0.5B run suggests it will not be answerable at that scale at all,
since the gate never opens there within 50 epochs.

\section{Discussion}
\label{sec:discussion}

\textbf{The knowing--using gap as a first-class object.} Our results
argue for treating the gap between internal recall and behavioral
production as a measurable training-time quantity rather than a
post-hoc evaluation surprise. The likelihood metric and the mid-layer
probe are cheap (a forward pass per epoch; a linear head), and together
they separate regimes that aggregate observables conflate:
not-yet-known, known-but-unsayable, and---in the $\lambda$-decay
arm---representation-known-but-output-never-formed.

\textbf{The gate signal and its schedule are design axes.} The same
formal gate behaves very differently depending on what it watches and
for how long. An output-likelihood gate defers engagement until the
model can recall---principled, but its onset is stochastic (one
replicate crossed at epoch 17; another never crossed), and a schedule
that silently never fires is indistinguishable from plain SFT. The
probe-saturation gate fires early at both scales because a fresh
linear head measures representation organization rather than LM-head
surface learning. What decided the outcome was not whether probe
gradients flowed, but \emph{when and for how long}: late,
output-triggered engagement at full weight matched SFT recall; early
sustained injection was high-variance and below SFT; and transient
injection that decayed to zero left the model at $0.012$ recall while
its mid layer organized the entities at $94.8\%$ decodability---even
though the arm then received twelve epochs of pure SFT and never
recovered. Probe regularization should therefore engage \emph{after}
saturation and then \emph{persist}; both halves of that prescription
are testable with the released harness.

\textbf{The saying gap is invariant to what we varied.} No arm at
either scale exceeded $1.2\%$ free-generation EM or $0.4\%$
generalization EM, while token-level recall under forcing reached
$99\%$: the gap lives at the representation-to-surface interface, and
closing it plausibly requires assembly-oriented objectives or decoding
interventions rather than better-organized storage.

\textbf{Controllers must not fight the phenomenon.} The decay-spiral
case study (Appendix~\ref{app:spiral}) is a small but concrete
instance of a general design error: an adaptive controller aimed at a
target regime treated that regime's signature as a fault and drove the
system away from it. The floor repair converts the controller from
opposed to subordinate: damp instability, never extinguish the
dynamics under study.

\section{Limitations}
\label{sec:limitations}

Our findings are bounded in four ways. First, the corpus is synthetic
and single-domain: crossing epochs and gap magnitudes will differ on
natural knowledge (e.g., long-tail Wikipedia facts
\citep{kandpal2023large}). Second, each schedule is represented by one
or two replicates, run-to-run variance at 1.7B is itself large
(sustained arms: $0.109$ vs.\ $0.805$), and GPU non-determinism makes
even same-seed runs diverge in takeoff epoch; our schedule claims rest
on paired tests against a common SFT control, but single-run
comparisons at this scale are unreliable. Third, the 24-epoch budget
censors late-takeoff arms: the never-engaged replicate was still
rising when it ended ($0.227$), so its endpoint is a lower bound.
Fourth, the probe is linear and trained on the corpus it monitors,
all experiments use LoRA at a single rank, and the mechanism behind
the decay-arm collapse---like that of the sustained arms'
variance---is hypothesized rather than established; the released
per-step telemetry is intended to make such follow-ups cheap.

\section{Conclusion}
\label{sec:conclusion}

We presented MLPR, a training-time instrumentation-and-control layer
that makes the knowing--using gap measurable online and acts on it
through a saturation gate. In a two-scale engagement-schedule grid we
showed that the gate signal decides when engagement is possible
(probe decodability fires at epochs 3--6, the output signal 11--16
epochs later or never), that the engagement schedule decides whether
memorization completes at all (final recall $0.012$--$0.934$ under
identical data, seed, and optimizer, with late output-triggered
engagement reaching SFT parity by McNemar test), captured the gate's
live re-evaluation transient, measured a knowing/saying dissociation
that no schedule closes (generalization EM $\le 0.4\%$ everywhere)
and a representation-without-output regime induced by transient
engagement, and confirmed from released adapters that memorization
updates concentrate in upper-layer MLP projections. Transferring the
instruments to natural corpora, scaling replicates beyond two per
arm, and pairing the regularized phase with assembly-oriented
objectives are immediate next steps.

\section*{Ethics Statement}
This work studies memorization dynamics, which sit at the intersection
of capability and privacy: the same teacher-forced recall metric that
monitors learning could, in adversarial hands, audit models for
memorized personal data, and the recall/production dissociation we
document cautions that exact-match audits may \emph{underestimate}
what a model has stored. We release only synthetic entities invented
for this study, and we believe better monitoring instruments
ultimately serve privacy auditing and data governance more than they
enable extraction.

\section*{Reproducibility Statement}
All code (dataset generator, trainer, callbacks, gate and engagement
schedules, controller, Modal deployment, analysis, McNemar module, and
figure scripts) is included in the released codebase, together with
the exact configurations of Appendix~\ref{app:hyper}. For every
production run we release the final adapter and probe head, the
per-epoch gate history, the complete stdout log, and a per-sample
final evaluation record whose prompt-aligned correctness vectors are
the direct inputs of every McNemar test reported here; per-epoch
numbers in every figure and in Appendix~\ref{app:curves} are
traceable to these artifacts, identified by their Weights \& Biases
IDs in Appendix~\ref{app:artifacts}. The earlier 50-epoch-class runs
behind the decay-spiral case study and the scale-saturation pilots
are the v4-generation artifacts also listed there.


\bibliography{refs}
\bibliographystyle{colm2026_conference}

\appendix
\section{Hyperparameters and Configurations}
\label{app:hyper}

Table~\ref{tab:hyper} lists the exact configurations used for both
production runs. Both jobs ran on a single NVIDIA H100 80~GB in bf16 via
the Modal cloud platform, with the full stdout stream teed to a
persistent volume and uploaded to the Hub after (and, on the 7B run,
during) training. The W\&B project \texttt{knowing\_using\_gap} holds
per-step and per-epoch curves, LoRA/probe matrix panels (every 5
epochs), and locked checkpoint artifacts.

\begin{table}[h]
\centering
\small
\begin{tabular*}{\columnwidth}{@{\extracolsep{\fill}} l l l}
\toprule
Parameter & 0.5B run & 7B run \\
\midrule
Model & Qwen2.5-0.5B-Instruct & Qwen2.5-7B-Instruct \\
LoRA $r$ / $\alpha$ / dropout & 16 / 32 / 0.05 & 16 / 32 / 0.05 \\
LoRA targets & q,k,v,o,gate,up,down & q,k,v,o,gate,up,down \\
Probe layer $l^\star$ & 12 of 24 & 14 of 28 \\
Probe head & $896 \to 1600$ & $3584 \to 1600$ \\
Backbone LR & $2{\times}10^{-5}$ & $2{\times}10^{-5}$ \\
Probe LR & $10^{-3}$ & $10^{-3}$ \\
Schedule & cosine, warmup $10\%$ & cosine, warmup $10\%$ \\
Batch (dev $\times$ accum) & $4 \times 4$ & $4 \times 4$ \\
Steps per epoch & 75 & 75 \\
Epochs (planned) & 50 & 50 \\
Epochs completed & 50 & 27 \\
Gate metric / $\tau_0$ & max / 0.8 & max / 0.8 \\
$\lambda_0$ / $\delta_0$ & 0.3 / 0.1 & 0.3 / 0.1 \\
Controller decay $\gamma$ & 0.7 & 0.7 \\
Damping floor $m_{\min}$ & 0.5 & 0.5 \\
Boost factor / cap & $1.5$ / $8\times$ & $1.5$ / $8\times$ \\
Lifecycle pushes & every 10 ep & every 10 ep \\
Matrices to W\&B & every 5 ep & every 5 ep \\
\bottomrule
\end{tabular*}
\caption{Full hyperparameter configurations.}
\label{tab:hyper}
\end{table}

\section{Dataset Construction}
\label{app:dataset}

The generator (released as
\texttt{scripts/generate\_synthetic\_dataset.py}) proceeds in three
steps. \textbf{(1) Entity vocabulary:} 1{,}600 invented entity names are
generated from syllable templates; each is assigned an id in
$[0, 1600)$ for the probe target space. \textbf{(2) Memorization split
$\mathcal{D}_{\mathrm{mem}}$:} 1{,}200 question--answer pairs in a fixed
interrogative template ending with \texttt{Answer:} \emph{label}; the
suffix is present during CE training (so the format is learned) and
stripped during all evaluations (so it cannot leak the answer into the
prompt). \textbf{(3) Generalization split $\mathcal{D}_{\mathrm{gen}}$}
(450 prompts): 200 \emph{chaining} prompts combine two trained facts
into a compositional question, and 250 \emph{intersection} prompts
re-ask trained facts with paraphrased templates and surface forms never
seen in training. Free-generation evaluations use left-padding,
greedy decoding, and strict exact match after answer extraction; the
likelihood evaluation \eqref{eq:lf} scores the first answer token and
the token-level overlap under teacher forcing. The per-epoch gate
evaluations run on a fixed random subsample of 208 memorization
prompts (13 batches) plus the full $\mathcal{D}_{\mathrm{gen}}$, keeping
epoch overhead at $\sim$1--2 minutes.

\section{Full Per-Epoch Curves}
\label{app:curves}

Table~\ref{tab:curves05} reports the 0.5B run every 5 epochs and
Table~\ref{tab:curves7} the 7B run every 3 epochs (the complete
per-epoch JSON histories are in the released telemetry).
Eval loss is the held-out generalization-split loss logged by the
trainer each epoch; $m(t)$ is the controller's cumulative LR
multiplier.

\begin{table}[h]
\centering
\scriptsize
\setlength{\tabcolsep}{3.2pt}
\begin{tabular*}{\columnwidth}{@{\extracolsep{\fill}} r rrrrrr}
\toprule
ep & $\amemlf$ & tok.\ & $\amemem$ & $\agen$ & $\lambda$ & eval \\
\midrule
1 & 0.000 & 0.337 & 0.000 & 0.000 & 0.000 & 5.03 \\
5 & 0.000 & 0.748 & 0.000 & 0.000 & 0.000 & 3.96 \\
10 & 0.010 & 0.758 & 0.000 & 0.000 & 0.000 & 4.23 \\
15 & 0.025 & 0.771 & 0.000 & 0.000 & 0.000 & 4.22 \\
20 & 0.030 & 0.786 & 0.005 & 0.000 & 0.000 & 4.40 \\
25 & 0.125 & 0.819 & 0.005 & 0.000 & 0.000 & 4.45 \\
30 & 0.240 & 0.857 & 0.055 & 0.000 & 0.000 & 4.50 \\
35 & 0.275 & 0.866 & 0.010 & 0.000 & 0.000 & 4.55 \\
40 & 0.355 & 0.895 & 0.050 & 0.000 & 0.000 & 4.57 \\
45 & 0.375 & 0.899 & 0.010 & 0.000 & 0.000 & 4.61 \\
50 & 0.380 & 0.900 & 0.025 & 0.000 & 0.000 & 4.64 \\
\bottomrule
\end{tabular*}
\caption{0.5B run (complete, 3{,}750 steps). $m(t)$: $1.0$ (ep 1--5),
$0.7$ (ep 6--7), $0.5$ (ep 8--50); one memorization diagnosis at the
controller's stagnation trigger. The gate never opens.}
\label{tab:curves05}
\end{table}

\begin{table}[h]
\centering
\scriptsize
\setlength{\tabcolsep}{3.2pt}
\begin{tabular*}{\columnwidth}{@{\extracolsep{\fill}} r rrrrrr}
\toprule
ep & $\amemlf$ & tok.\ & $\amemem$ & $\agen$ & $\lambda$ & eval \\
\midrule
1 & 0.000 & 0.386 & 0.000 & 0.000 & 0.000 & 5.10 \\
3 & 0.015 & 0.752 & 0.000 & 0.000 & 0.000 & 3.21 \\
6 & 0.005 & 0.760 & 0.000 & 0.000 & 0.000 & 4.02 \\
9 & 0.025 & 0.783 & 0.000 & 0.000 & 0.000 & 4.06 \\
12 & 0.110 & 0.822 & 0.005 & 0.000 & 0.000 & 4.19 \\
15 & 0.320 & 0.895 & 0.015 & 0.000 & 0.000 & 4.24 \\
18 & 0.570 & 0.946 & 0.030 & 0.000 & 0.000 & 4.34 \\
21 & 0.745 & 0.971 & 0.030 & 0.000 & 0.000 & 4.40 \\
24 & \textbf{0.805} & 0.979 & 0.055 & 0.000 & 0.015 & 4.44 \\
25 & 0.825 & 0.981 & 0.025 & 0.000 & 0.075 & 4.41 \\
27 & 0.875 & 0.987 & 0.020 & 0.000 & 0.225 & --- \\
\bottomrule
\end{tabular*}
\caption{7B run (27 completed epochs; terminated during epoch 28 by
compute-budget exhaustion, hence no epoch-28 eval). The crossing of
$\tau_0{=}0.8$ at epoch 24 opens the gate; $\lambda(t)$ ramps to
$0.225$ by epoch 27. Epoch-27 values were read from the live run log.}
\label{tab:curves7}
\end{table}

\section{The Decay-Spiral Failure and Its Repair}
\label{app:spiral}

\textbf{Failure.} In a pilot 50-epoch 0.5B run (v3), the eval-aware
controller issued eval-driven decays every two epochs for five
consecutive episodes while held-out loss diverged, driving the
cumulative LR multiplier along
$1.0 \to 0.49 \to 0.34 \to 0.24 \to 0.17$, breaching the controller's
own nominal minimum ($0.3$) because the eval-driven pathway applied
decays without consulting the instability path's floor. Memorization
was choked: likelihood recall stayed pinned near $0.02$ and the gate
threshold became unreachable. The run was stopped to avoid burning
compute on dynamics that the controller itself had frozen.

\textbf{Diagnosis.} The failure is structural, not incidental: in the
memorization regime, held-out divergence is the \emph{expected}
signature that the saturation gate waits for, so an eval-driven
damping rule without a floor is guaranteed to eventually fight the
gate.

\textbf{Repair.} The eval-driven pathway now clamps the cumulative
multiplier to $m_{\min}{=}0.5$ (Eq.~\ref{eq:floor}) and refuses
further eval-driven decays at the floor. Validation proceeded in
three stages: (i) unit tests on a stubbed trainer assert the floor is
respected exactly ($m$ lands on $0.5$, never below; optimizer LRs
match base $\times$ floor) among 138 checks; (ii) a dedicated 5-epoch
GPU run with an unreachable gate reproduced the failing regime
($\lambda{=}0$ throughout, eval rising) and showed one eval-driven
decay ($1.0 \to 0.7$) with the floor never breached; (iii) both
production runs in this paper clamp at $0.5$ and hold
(Fig.~\ref{fig:dyn}c) while memorization proceeds unimpeded to
saturation at 7B.

\section{Artifact Map and Log Reconstruction}
\label{app:artifacts}

\textbf{Hugging Face Hub (user \texttt{codewithdark}).}
\texttt{mlpr-qwen2.5-0.5b-instruct-50ep-adaptive-v4}: final adapter,
probe head (\texttt{probe\_head.pt}), tokenizer, and full telemetry
(\texttt{gate\_history.json}, \texttt{adaptive\_history.json},
\texttt{trainer\_state.json} with 1{,}925 log entries,
\texttt{training\_log.txt}, 1.5\,MB). The 7B repository
\texttt{mlpr-qwen2.5-7b-instruct-50ep-adaptive-v4} holds the epoch-26
adapter (the trainer's last in-progress hub commit) and the training
log through epoch 26; its per-epoch gate and controller records for
epochs 1--25 are parsed from the log's structured
\texttt{[MEM\ GATE]} / \texttt{[ADAPTIVE]} lines, and the epoch-27
record was captured verbatim from the live log stream before
termination. \textbf{Weights \& Biases} (project
\texttt{knowing\_using\_gap}): runs
\texttt{mlpr-qwen2.5-0.5b-50ep-adaptive-v4} and
\texttt{mlpr-qwen2.5-7b-50ep-adaptive-v4} with per-step losses,
learning rates, gradient norms, per-epoch adaptive decisions, LoRA/probe
matrix panels every 5 epochs (168 heatmaps each of $B$, $A$, and
$\Delta W$), and locked checkpoint artifacts.
\textbf{Codebase}: the released zip contains the full training stack
(\texttt{src/} with probe, LoRA setup, gate callback, adaptive
controller, likelihood evaluator, W\&B matrix logger), the Modal
deployment (\texttt{modal\_train.py}), the dataset generator, the
analysis and figure scripts, and this paper under \texttt{paper/}.

\section{The Likelihood Recall Metric in Detail}
\label{app:lfmetric}

For a memorization pair $(x, y)$ the prompt is built by stripping the
\texttt{Answer:} suffix from the training text. The metric performs a
single teacher-forced forward pass over
$\texttt{prompt} \oplus \texttt{answer}$ with left padding and records
(i) strict first-answer-token exact match \eqref{eq:lf}, and (ii) the
fraction of answer tokens whose greedy next-token prediction matches
the reference (token recall). Both are macro-averaged over the
subsample. Compared with free-generation EM, this removes decoding
length, format, and repetition effects; the residual gap between token
recall and strict EM isolates single-token identity from multi-token
assembly. Substring containment of the target in the free generation
complements both by detecting \emph{stored but misframed} answers---the
$76.5\%$ vs.\ $2\%$ contrast in Table~\ref{tab:main}.


\end{document}
